\chapter{Skema Bintang Pertama}
\label{modul:skema-bintang-pertama}

\section{Identitas Modul}

\begin{identitasmodul}
\textbf{Sumber materi}    & Pertemuan 3 \\
\textbf{CPMK}             & CPMK-072 \\
\textbf{Minggu pelaksanaan} & Minggu ke-4 \\
\textbf{Durasi tatap muka} & 120 menit \\
\textbf{Estimasi tugas}   & 3--4 jam di luar sesi \\
\textbf{Moda kerja}       & Individual \\
\textbf{Perangkat}        & PostgreSQL 17 dan DBeaver \\
\textbf{Dataset}          & NusaMart ukuran kecil \\
\textbf{Skrip pendukung}  & \berkas{sql/modul-02/} beserta \berkas{check.sql} \\
\textbf{Luaran}           & \berkas{ddl.sql}, \berkas{load.sql}, dan \berkas{query.sql} \\
\end{identitasmodul}

\slideref{3}

\section{Capaian dan Indikator Keberhasilan}

Mahasiswa mampu menerjemahkan satu proses bisnis menjadi skema bintang,
membangun dimensi tanggal, memuat dimensi dan fakta, serta merekonsiliasi hasil
star query dengan sumber. Keberhasilan ditandai oleh hasil agregat yang sama
persis dengan basis data sumber.

\begin{capaian}
Pada akhir modul ini mahasiswa mampu:
\begin{enumerate}[leftmargin=1.4em]
  \item membangkitkan \perintah{dim\_tanggal} beserta atribut kalender yang
        diperlukan analisis, bukan sekadar kolom tanggal;
  \item membedakan \istilah{natural key} dari \istilah{surrogate key} dan
        menjelaskan mengapa tabel fakta menunjuk surrogate key;
  \item menulis DDL dimensi dan fakta yang grain-nya sesuai dengan pernyataan
        grain Modul 1;
  \item memuat dimensi dan fakta dari sumber dengan pemetaan kunci yang benar;
  \item menulis star query dan membuktikan angkanya sama persis dengan query
        ekuivalen pada basis data sumber.
\end{enumerate}
\end{capaian}

Rekonsiliasi adalah inti modul ini. Skema bintang yang indah tetapi
menghasilkan angka berbeda dari sumber tidak bernilai apa pun --- dan selisih
itu hampir selalu berasal dari satu keputusan yang tidak dicatat, bukan dari
kesalahan SQL.

\section{Prasyarat dan Persiapan}

\subsection{Prasyarat}

\begin{itemize}
  \item Pernyataan grain, kandidat measure, dan kandidat dimensi hasil Modul 1;
  \item DDL PostgreSQL: \perintah{CREATE TABLE}, tipe data dasar, primary key,
        dan foreign key;
  \item agregasi SQL dengan \perintah{GROUP BY} beserta \perintah{JOIN} banyak
        tabel;
  \item pemahaman bahwa satu baris pada \perintah{transaksi\_item} adalah satu
        produk pada satu struk --- temuan Bagian C Modul 1.
\end{itemize}

\subsection{Pre-lab}

\begin{enumerate}[leftmargin=1.4em]
  \item Pastikan \berkas{desain-konseptual-v1.md} sudah diperiksa asisten. Grain
        yang salah pada modul ini akan terbawa sampai Modul 10.
  \item Bila Modul 1 belum selesai, pulihkan \berkas{snapshot-modul-01.sql}
        agar basis data sumber berada pada keadaan yang sama dengan kelas.
  \item Pastikan basis data \perintah{nusamart\_dw} dapat diakses dan masih
        kosong.
  \item Siapkan jawaban atas pertanyaan berikut di dalam
        \berkas{modul-02-star-schema/pre-lab.md}:
  \begin{enumerate}[label=\alph*.,leftmargin=1.4em]
    \item Mengapa tabel fakta tidak menyimpan nama produk secara langsung?
    \item Berapa baris yang Anda harapkan ada pada \perintah{fakta\_penjualan},
          dan dari query mana angka itu berasal?
    \item Transaksi berstatus batal akan Anda muat atau saring? Tuliskan
          keputusannya sekarang, sebelum menulis DDL.
  \end{enumerate}
\end{enumerate}

\section{Konsep Dasar}

\subsection{Anatomi Skema Bintang}

Skema bintang menempatkan satu tabel fakta di pusat, dikelilingi tabel dimensi
yang menunjuk langsung ke sana tanpa tabel perantara \citep{kimball2013}. Bentuk
ini lahir dari satu pertimbangan praktis: setiap join tambahan adalah biaya, dan
query analitik menjalankan join yang sama berulang kali.

\begin{konsep}{Pembagian tugas fakta dan dimensi}
\begin{itemize}
  \item \textbf{Tabel fakta} menyimpan \emph{apa yang terjadi} --- besaran
        numerik yang diukur, ditambah kunci asing ke setiap dimensi. Barisnya
        banyak, kolomnya sedikit, dan hampir seluruh kolomnya berupa angka.
  \item \textbf{Tabel dimensi} menyimpan \emph{konteks} --- atribut yang dipakai
        untuk menyaring, mengelompokkan, dan memberi label. Barisnya sedikit,
        kolomnya banyak, dan sebagian besar kolomnya berupa teks.
\end{itemize}
Aturan praktis: kolom yang muncul setelah \perintah{SELECT} dan
\perintah{GROUP BY} berasal dari dimensi; kolom di dalam \perintah{SUM} berasal
dari fakta.
\end{konsep}

Dimensi sengaja \emph{tidak} dinormalisasi. Pada sumber, kategori produk
tersimpan pada tabel \perintah{kategori} yang mengacu pada dirinya sendiri, dan
mengambil nama kategori memerlukan join berjenjang. Pada gudang data, hierarki
itu diratakan menjadi kolom \perintah{kategori} dan \perintah{sub\_kategori} di
dalam \perintah{dim\_produk}. Redundansi yang timbul memang disengaja: ia
ditukar dengan hilangnya join, dan pada tabel berisi 180 ribu baris biayanya
tidak berarti.

\begin{catatan}
Bentuk yang mempertahankan normalisasi dimensi disebut \istilah{snowflake}.
Ia menghemat ruang tetapi menambah join, dan pada gudang data ruang jauh lebih
murah daripada waktu. Modul ini memakai skema bintang; bila Anda tergoda
memecah dimensi menjadi beberapa tabel, tuliskan dulu berapa join tambahan yang
harus dibayar setiap query.
\end{catatan}

\subsection{Natural Key dan Surrogate Key}

\istilah{Natural key} adalah kunci yang dipakai sistem operasional ---
\perintah{produk\_id} pada tabel \perintah{produk}. \istilah{Surrogate key}
adalah kunci buatan gudang data, berupa bilangan bulat tanpa makna bisnis,
misalnya \perintah{produk\_key}.

Tabel fakta menunjuk surrogate key, bukan natural key. Ada tiga alasan, dan
alasan ketiga adalah yang paling menentukan:

\begin{enumerate}[leftmargin=1.4em]
  \item \textbf{Ukuran.} Bilangan bulat lebih sempit daripada kode teks, dan
        tabel fakta memuat jutaan baris.
  \item \textbf{Kekebalan terhadap sumber.} Bila NusaMart mengganti sistem
        kasir dan seluruh kode produk berubah, gudang data tidak perlu ditulis
        ulang.
  \item \textbf{Kemampuan menyimpan histori.} Satu produk dapat memiliki
        beberapa versi ketika harganya berubah. Ketiga versi memiliki
        \perintah{produk\_id} yang sama tetapi \perintah{produk\_key} yang
        berbeda --- dan hanya karena itulah tabel fakta dapat menunjuk versi
        yang berlaku saat transaksi terjadi.
\end{enumerate}

\begin{catatan}
Alasan ketiga baru akan terasa manfaatnya pada Modul 3, ketika
\perintah{dim\_produk} diubah menjadi SCD Type 2. Pilihan yang Anda ambil hari
ini menyiapkan jalan untuk itu. Karena itu, jangan pernah memakai
\perintah{produk\_id} sebagai primary key \perintah{dim\_produk}, meskipun hari
ini tampak lebih sederhana.
\end{catatan}

Natural key tetap disimpan di dalam dimensi sebagai kolom biasa. Ia diperlukan
saat pemuatan --- untuk mencari surrogate key yang sesuai --- dan saat
menelusuri satu angka kembali ke sumbernya.

\subsection{Dimensi Tanggal}

Dimensi tanggal adalah satu-satunya dimensi yang dibangkitkan, bukan diambil
dari sumber. Alasannya sederhana: tidak ada sistem operasional yang menyimpan
tabel kalender, padahal hampir setiap pertanyaan analitik memakai waktu.

Menyimpan tanggal sebagai kolom \perintah{DATE} di tabel fakta tampak cukup,
tetapi tidak. Pertanyaan seperti ``bagaimana penjualan akhir pekan dibandingkan
hari kerja?'' memerlukan atribut yang harus dihitung ulang pada setiap query
bila kalendernya tidak ada.

\begin{konsep}{Atribut kalender minimum}
\begin{itemize}
  \item Kunci: \perintah{tanggal\_key} dalam bentuk \perintah{YYYYMMDD}, dan
        \perintah{tanggal} sebagai nilai \perintah{DATE} sebenarnya.
  \item Penanggalan: \perintah{hari}, \perintah{nama\_hari}, \perintah{bulan},
        \perintah{nama\_bulan}, \perintah{kuartal}, \perintah{tahun}.
  \item Pengelompokan: \perintah{tahun\_bulan} untuk pengurutan,
        \perintah{minggu\_ke}, dan \perintah{hari\_ke\_tahun}.
  \item Penanda: \perintah{akhir\_pekan} bertipe \perintah{BOOLEAN}.
\end{itemize}
\end{konsep}

Kunci \perintah{YYYYMMDD} adalah pengecualian yang diterima atas aturan
``surrogate key tanpa makna''. Bentuk ini memudahkan pembacaan saat
\istilah{debugging} dan memungkinkan partisi berdasarkan rentang kunci pada
Modul 6, tanpa kehilangan sifatnya sebagai bilangan bulat sempit.

\subsection{Measure dan Additivity}

\istilah{Measure} adalah kolom numerik pada tabel fakta. Sifat penting yang
harus diketahui sebelum menyimpannya adalah \istilah{additivity} --- terhadap
dimensi apa saja measure itu boleh dijumlahkan.

\begin{center}
\small
\begin{tabular}{@{}p{0.20\textwidth}>{\raggedright\arraybackslash}p{0.34\textwidth}>{\raggedright\arraybackslash}p{0.38\textwidth}@{}}
\toprule
\textbf{Sifat} & \textbf{Arti} & \textbf{Contoh pada NusaMart} \\
\midrule
Additive & Boleh dijumlahkan terhadap seluruh dimensi &
  \perintah{kuantitas}, \perintah{subtotal}, \perintah{diskon} \\
Semi-additive & Boleh dijumlahkan terhadap sebagian dimensi, tetapi tidak
  terhadap waktu & Saldo persediaan --- dibahas pada Modul 4 \\
Non-additive & Tidak boleh dijumlahkan sama sekali &
  \perintah{harga\_satuan}, rasio, dan persentase \\
\bottomrule
\end{tabular}
\end{center}

Measure non-additive tetap disimpan. Yang keliru bukan menyimpannya, melainkan
menjumlahkannya. Harga satuan diperlukan untuk menghitung harga rata-rata
tertimbang, dan rata-rata tertimbang dihitung dari dua measure additive:
$\sum(\text{subtotal}) / \sum(\text{kuantitas})$.

\begin{peringatan}
Rasio tidak pernah disimpan sebagai measure. Rata-rata dari rata-rata bukanlah
rata-rata: menjumlahkan sepuluh nilai ``margin persen'' lalu membaginya dengan
sepuluh menghasilkan angka yang salah kecuali seluruh transaksinya berukuran
sama. Simpan pembilang dan penyebutnya sebagai measure additive, lalu hitung
rasionya saat query.
\end{peringatan}

\section{Studi Kasus dan Protokol Praktikum}

Proses penjualan NusaMart digunakan untuk membangun skema bintang pertama.
Seluruh query analitik wajib dibandingkan dengan query ekuivalen pada sumber.

Skema yang dibangun mengikuti pernyataan grain Modul 1:

\begin{quote}
\emph{Satu baris fakta mewakili satu produk pada satu baris struk penjualan di
satu toko pada satu waktu.}
\end{quote}

Grain ini berpadanan satu lawan satu dengan baris pada
\perintah{nusamart.transaksi\_item}. Karena itu jumlah baris
\perintah{fakta\_penjualan} harus sama dengan jumlah baris
\perintah{transaksi\_item} yang lolos penyaringan status --- dan angka itu
adalah pemeriksaan rekonsiliasi pertama Anda.

\begin{center}
\small
\begin{tabular}{@{}p{0.26\textwidth}>{\raggedright\arraybackslash}p{0.30\textwidth}>{\raggedright\arraybackslash}p{0.36\textwidth}@{}}
\toprule
\textbf{Tabel gudang} & \textbf{Sumber} & \textbf{Catatan} \\
\midrule
\perintah{dim\_tanggal} & dibangkitkan &
  Sepuluh tahun, 2020 sampai 2029 \\
\perintah{dim\_produk} & \perintah{produk} dan \perintah{kategori} &
  Hierarki kategori diratakan \\
\perintah{dim\_toko} & \perintah{toko} &
  Diambil apa adanya \\
\perintah{fakta\_penjualan} & \perintah{transaksi\_item} dan
  \perintah{transaksi} & Grain mengikuti \perintah{transaksi\_item} \\
\bottomrule
\end{tabular}
\end{center}

Dimensi pelanggan sengaja \emph{belum} dibangun. Dari Modul 1 diketahui bahwa
sebagian besar transaksi tidak memiliki \perintah{pelanggan\_id}, dan penanganan
yang benar memerlukan baris \perintah{unknown} yang baru dibahas pada Modul 5.
Menundanya bukan penyederhanaan yang malas, melainkan urutan yang disengaja.

\begin{catatanasisten}
Seluruh objek modul ini dibangun pada schema \perintah{public} basis data
\perintah{nusamart\_dw}. Penataan ke schema \perintah{staging},
\perintah{gudang}, dan \perintah{mart} adalah pekerjaan Modul 5; bila
dikerjakan sekarang, Modul 5 kehilangan bahan latihannya.
\end{catatanasisten}

\section{Alur Praktikum 120 Menit}

\begin{alurwaktu}
0--15 & Demonstrasi grain dan bentuk skema bintang. \\
15--95 & Kerja terbimbing: DDL, pemuatan, dan star query. \\
95--110 & Checkpoint struktur tabel dan rekonsiliasi hasil. \\
110--120 & Penjelasan tugas proses bisnis kedua. \\
\end{alurwaktu}

\section{Prosedur Praktikum Terbimbing}

Seluruh perintah pada bagian ini dijalankan pada \perintah{nusamart\_dw}, kecuali
disebutkan lain. Kumpulkan DDL pada \berkas{ddl.sql}, perintah pemuatan pada
\berkas{load.sql}, dan query analitik pada \berkas{query.sql} --- ketiganya harus
dapat dijalankan berurutan dari awal.

\subsection{Bagian A: Membentuk Dimensi Tanggal}

\langkah{A.1 Membangkitkan kalender sepuluh tahun}{10 menit}

Kalender sepuluh tahun dibangkitkan dengan \perintah{generate\_series} oleh
skrip \berkas{sql/\allowbreak modul-02/\allowbreak 00-dim-tanggal.sql}.
Jalankan skrip itu, lalu baca isinya --- Anda perlu memahami cara kerjanya,
bukan sekadar menjalankannya.

\begin{lstlisting}[style=sqlgaya]
CREATE TABLE dim_tanggal (
  tanggal_key    INTEGER     PRIMARY KEY,
  tanggal        DATE        NOT NULL UNIQUE,
  hari           SMALLINT    NOT NULL,
  nama_hari      VARCHAR(10) NOT NULL,
  minggu_ke      SMALLINT    NOT NULL,
  bulan          SMALLINT    NOT NULL,
  nama_bulan     VARCHAR(12) NOT NULL,
  kuartal        SMALLINT    NOT NULL,
  tahun          SMALLINT    NOT NULL,
  tahun_bulan    INTEGER     NOT NULL,
  hari_ke_tahun  SMALLINT    NOT NULL,
  akhir_pekan    BOOLEAN     NOT NULL
);

INSERT INTO dim_tanggal
SELECT TO_CHAR(d, 'YYYYMMDD')::INTEGER,
       d::DATE,
       EXTRACT(DAY     FROM d),
       TO_CHAR(d, 'TMDay'),
       EXTRACT(WEEK    FROM d),
       EXTRACT(MONTH   FROM d),
       TO_CHAR(d, 'TMMonth'),
       EXTRACT(QUARTER FROM d),
       EXTRACT(YEAR    FROM d),
       TO_CHAR(d, 'YYYYMM')::INTEGER,
       EXTRACT(DOY     FROM d),
       EXTRACT(ISODOW  FROM d) IN (6, 7)
FROM   generate_series(DATE '2020-01-01', DATE '2029-12-31', INTERVAL '1 day') d;
\end{lstlisting}

\langkah{A.2 Memverifikasi kalender}{5 menit}

\begin{lstlisting}[style=sqlgaya]
SELECT COUNT(*)          AS jumlah_hari,
       MIN(tanggal)      AS mulai,
       MAX(tanggal)      AS selesai,
       COUNT(*) FILTER (WHERE akhir_pekan) AS jumlah_akhir_pekan
FROM   dim_tanggal;
\end{lstlisting}

Sepuluh tahun 2020--2029 memuat 3.653 hari karena 2020, 2024, dan 2028 adalah
tahun kabisat. Bila angka Anda 3.650, rentangnya keliru.

\begin{peringatan}
Penanda \perintah{TM} pada \perintah{TO\_CHAR} menghasilkan nama hari dan bulan
mengikuti \perintah{lc\_time} basis data. Bila keluarannya berbahasa Inggris
padahal Anda mengharapkan bahasa Indonesia, itu bukan kesalahan skrip melainkan
setelan lokal. Yang penting adalah konsistensi: pilih satu, lalu jangan
mencampurnya di tengah semester.
\end{peringatan}

\subsection{Bagian B: Menyalin Sumber dan Membuat Tabel Dimensi}

\langkah{B.1 Menyalin sumber ke tabel staging}{10 menit}

Sumber dan gudang berada pada dua basis data terpisah, sehingga isinya harus
dipindahkan lebih dahulu. Skrip \berkas{sql/modul-02/01-staging.sql} membuat
empat tabel penampung pada \perintah{nusamart\_dw}:
\perintah{staging\_produk}, \perintah{staging\_kategori},
\perintah{staging\_toko}, dan \perintah{staging\_penjualan}.

Tabel staging menyalin sumber \textbf{apa adanya} --- tanpa surrogate key,
tanpa perbaikan, dan tanpa penyaringan. Pembersihan dan pemetaan kunci terjadi
pada langkah berikutnya, dan pemisahan ini yang membuat pemuatan dapat diulang
tanpa menyentuh sumber kembali.

\begin{lstlisting}[style=shellgaya]
# Ekspor dari sumber, lalu impor ke gudang. Ulangi untuk keempat tabel.
docker compose exec -T postgres_source \
  psql -U praktikum -d nusamart_oltp \
  -c "COPY (SELECT * FROM nusamart.toko) TO STDOUT WITH CSV" \
| docker compose exec -T postgres_dw \
  psql -U praktikum -d nusamart_dw \
  -c "COPY staging_toko FROM STDIN WITH CSV"
\end{lstlisting}

\begin{catatan}
Perintah \perintah{COPY ... TO STDOUT} yang disalurkan langsung ke
\perintah{COPY ... FROM STDIN} memindahkan data tanpa singgah sebagai berkas.
Opsi \perintah{-T} diperlukan agar Docker tidak mengalokasikan terminal, yang
akan merusak aliran data biner.
\end{catatan}

\langkah{B.2 Menulis DDL dimensi}{10 menit}

Perhatikan pemisahan surrogate key dan natural key pada kedua tabel berikut.

\begin{lstlisting}[style=sqlgaya]
CREATE TABLE dim_produk (
  produk_key    INTEGER GENERATED ALWAYS AS IDENTITY PRIMARY KEY,
  produk_id     INTEGER      NOT NULL UNIQUE,   -- natural key dari sumber
  sku           VARCHAR(20)  NOT NULL,
  nama_produk   VARCHAR(150) NOT NULL,
  merek         VARCHAR(60),
  kategori      VARCHAR(60),                    -- hierarki yang diratakan
  sub_kategori  VARCHAR(60),
  satuan        VARCHAR(15),
  berat_gram    NUMERIC(10,2)
);

CREATE TABLE dim_toko (
  toko_key      INTEGER GENERATED ALWAYS AS IDENTITY PRIMARY KEY,
  toko_id       INTEGER     NOT NULL UNIQUE,
  kode_toko     VARCHAR(15) NOT NULL,
  nama_toko     VARCHAR(100) NOT NULL,
  kota          VARCHAR(60),
  provinsi      VARCHAR(60),
  tipe_toko     VARCHAR(30),
  tanggal_buka  DATE
);
\end{lstlisting}

\langkah{B.3 Memuat dimensi dari staging}{10 menit}

Dimensi diisi dari staging, bukan dari sumber secara langsung. Di sinilah
hierarki kategori diratakan menjadi dua kolom.

\begin{lstlisting}[style=sqlgaya]
INSERT INTO dim_produk (produk_id, sku, nama_produk, merek,
                        kategori, sub_kategori, satuan, berat_gram)
SELECT     p.produk_id, p.sku, p.nama_produk, p.merek,
           induk.nama_kategori AS kategori,
           k.nama_kategori     AS sub_kategori,
           p.satuan, p.berat_gram
FROM       staging_produk p
LEFT JOIN  staging_kategori k     ON k.kategori_id     = p.kategori_id
LEFT JOIN  staging_kategori induk ON induk.kategori_id = k.induk_kategori_id;
\end{lstlisting}

Perhatikan join ganda ke \perintah{staging\_kategori}: inilah bentuk konkret
dari perataan hierarki. Satu join mengambil sub-kategori, satu lagi mengambil
induknya. Pada sumber, keduanya berada di satu tabel yang mengacu pada dirinya
sendiri; pada dimensi, keduanya menjadi dua kolom yang siap dipakai
\perintah{GROUP BY}.

\begin{peringatan}
Dari Modul 1 diketahui ada produk yang menunjuk kategori tidak dikenal.
\perintah{LEFT JOIN} membuat baris itu tetap masuk dengan kategori
\perintah{NULL}, sedangkan \perintah{INNER JOIN} akan membuangnya diam-diam ---
dan rekonsiliasi Anda akan meleset tanpa jejak. Gunakan \perintah{LEFT JOIN},
lalu hitung berapa baris yang berkategori \perintah{NULL} dan catat angkanya.
\end{peringatan}

\subsection{Bagian C: Membuat dan Memuat Fakta Penjualan}

\langkah{C.1 Menulis DDL fakta}{10 menit}

Tulis pernyataan grain sebagai komentar di baris pertama DDL. Ini bukan
formalitas: DDL yang tidak menyebutkan grainnya akan disalahpahami oleh orang
berikutnya, termasuk oleh Anda sendiri tiga modul kemudian.

\begin{lstlisting}[style=sqlgaya]
-- Grain: satu baris mewakili satu produk pada satu baris struk penjualan
--        di satu toko pada satu waktu.
CREATE TABLE fakta_penjualan (
  tanggal_key    INTEGER      NOT NULL REFERENCES dim_tanggal(tanggal_key),
  produk_key     INTEGER      NOT NULL REFERENCES dim_produk(produk_key),
  toko_key       INTEGER      NOT NULL REFERENCES dim_toko(toko_key),
  nomor_struk    VARCHAR(25)  NOT NULL,   -- degenerate dimension
  transaksi_id   BIGINT       NOT NULL,   -- penelusuran balik ke sumber
  kuantitas      NUMERIC(10,2) NOT NULL,
  harga_satuan   NUMERIC(12,2) NOT NULL,  -- non-additive
  diskon         NUMERIC(12,2) NOT NULL DEFAULT 0,
  subtotal       NUMERIC(14,2) NOT NULL
);
\end{lstlisting}

Kolom \perintah{nomor\_struk} adalah \istilah{degenerate dimension}: atribut
dimensi yang tidak memiliki tabel dimensi sendiri karena tidak punya atribut
lain untuk disimpan. Ia tetap berguna --- tanpanya, pertanyaan ``berapa
rata-rata jumlah produk per struk?'' tidak dapat dijawab.

\begin{catatan}
Dari Modul 1 diketahui \perintah{nomor\_struk} hanya unik pada pasangan
\perintah{(toko\_id, nomor\_struk)}. Karena \perintah{toko\_key} sudah ada pada
tabel fakta, pasangan \perintah{(toko\_key, nomor\_struk)} tetap
mengidentifikasi satu struk secara unik. Inilah contoh temuan Modul 1 yang
langsung memengaruhi rancangan Modul 2.
\end{catatan}

\langkah{C.2 Memuat fakta dengan pemetaan kunci}{15 menit}

Pemuatan fakta bukan penyalinan. Setiap natural key dari sumber harus
diterjemahkan menjadi surrogate key melalui join ke dimensi.

\begin{lstlisting}[style=sqlgaya]
INSERT INTO fakta_penjualan (
  tanggal_key, produk_key, toko_key, nomor_struk, transaksi_id,
  kuantitas, harga_satuan, diskon, subtotal)
SELECT dt.tanggal_key,
       dp.produk_key,
       dtk.toko_key,
       s.nomor_struk,
       s.transaksi_id,
       s.kuantitas,
       s.harga_satuan,
       s.diskon,
       s.subtotal
FROM   staging_penjualan s
JOIN   dim_tanggal dt  ON dt.tanggal = s.waktu_transaksi::DATE
JOIN   dim_produk  dp  ON dp.produk_id = s.produk_id
JOIN   dim_toko    dtk ON dtk.toko_id  = s.toko_id
WHERE  s.status <> 'BATAL';
\end{lstlisting}

\begin{peringatan}
\perintah{JOIN} biasa pada perintah di atas akan \emph{membuang} baris yang
tidak menemukan pasangan dimensi, tanpa pesan galat apa pun. Karena itu
langkah berikutnya wajib: hitung selisih antara jumlah baris sumber dan jumlah
baris yang berhasil dimuat. Bila selisihnya bukan nol, cari tahu sebabnya
sebelum melanjutkan. Penanganan yang benar --- mengarahkan baris tanpa pasangan
ke baris \perintah{unknown} --- dibahas pada Modul 5.
\end{peringatan}

\langkah{C.3 Memeriksa jumlah baris}{5 menit}

\begin{lstlisting}[style=sqlgaya]
SELECT (SELECT COUNT(*) FROM staging_penjualan WHERE status <> 'BATAL')
         AS baris_sumber,
       (SELECT COUNT(*) FROM fakta_penjualan) AS baris_fakta,
       (SELECT COUNT(*) FROM staging_penjualan WHERE status <> 'BATAL')
       - (SELECT COUNT(*) FROM fakta_penjualan) AS selisih;
\end{lstlisting}

Catat angka \perintah{selisih} pada laporan meskipun nilainya nol. Selisih yang
tidak dijelaskan adalah kegagalan; selisih yang dijelaskan adalah temuan.

\subsection{Bagian D: Menulis dan Merekonsiliasi Star Query}

\langkah{D.1 Menulis tiga star query}{15 menit}

Tulis tiga query yang menjawab pertanyaan bisnis berikut. Ketiganya harus
memakai bentuk baku skema bintang: fakta di-join ke dimensi, dikelompokkan
menurut atribut dimensi, dan diagregasi atas measure fakta.

\begin{enumerate}[leftmargin=1.4em]
  \item Total nilai penjualan per kategori produk per bulan sepanjang 2024.
  \item Sepuluh produk dengan nilai penjualan tertinggi di Provinsi Lampung.
  \item Perbandingan nilai penjualan hari kerja dan akhir pekan per tipe toko.
\end{enumerate}

\begin{lstlisting}[style=sqlgaya]
-- Contoh bentuk baku untuk pertanyaan pertama.
SELECT dp.kategori,
       dt.tahun_bulan,
       SUM(f.subtotal)  AS nilai_penjualan,
       SUM(f.kuantitas) AS unit_terjual
FROM   fakta_penjualan f
JOIN   dim_produk  dp ON dp.produk_key  = f.produk_key
JOIN   dim_tanggal dt ON dt.tanggal_key = f.tanggal_key
WHERE  dt.tahun = 2024
GROUP  BY dp.kategori, dt.tahun_bulan
ORDER  BY dp.kategori, dt.tahun_bulan;
\end{lstlisting}

Perhatikan bahwa query ini tidak memuat satu pun subquery dan hanya
memerlukan dua join. Query setara pada basis data sumber memerlukan empat join
dan satu join berjenjang untuk mencapai nama kategori. Selisih itulah yang
dibeli oleh denormalisasi.

\langkah{D.2 Merekonsiliasi dengan sumber}{15 menit}

Untuk setiap star query, tulis query ekuivalen pada \perintah{nusamart\_oltp},
lalu bandingkan hasilnya baris demi baris. Angkanya harus sama persis, bukan
sekadar mirip.

\begin{lstlisting}[style=sqlgaya]
-- Query pembanding, dijalankan pada nusamart_oltp.
SELECT induk.nama_kategori AS kategori,
       TO_CHAR(t.waktu_transaksi, 'YYYYMM')::INTEGER AS tahun_bulan,
       SUM(ti.subtotal)  AS nilai_penjualan,
       SUM(ti.kuantitas) AS unit_terjual
FROM       nusamart.transaksi_item ti
JOIN       nusamart.transaksi t     ON t.transaksi_id = ti.transaksi_id
JOIN       nusamart.produk p        ON p.produk_id    = ti.produk_id
LEFT JOIN  nusamart.kategori k      ON k.kategori_id  = p.kategori_id
LEFT JOIN  nusamart.kategori induk  ON induk.kategori_id = k.induk_kategori_id
WHERE      t.status <> 'BATAL'
  AND      EXTRACT(YEAR FROM t.waktu_transaksi) = 2024
GROUP  BY induk.nama_kategori, TO_CHAR(t.waktu_transaksi, 'YYYYMM')
ORDER  BY 1, 2;
\end{lstlisting}

\begin{catatan}
Rekonsiliasi yang meleset hampir selalu berasal dari tiga sebab, berurutan
menurut seringnya: penyaringan status yang berbeda antara kedua sisi,
\perintah{INNER JOIN} yang membuang baris pada salah satu sisi, dan pembulatan
\perintah{NUMERIC} yang berbeda. Periksa ketiganya dengan urutan itu sebelum
mencurigai hal lain.
\end{catatan}

\begin{luaran}
\berkas{modul-02-star-schema/ddl.sql}, \berkas{load.sql}, \berkas{query.sql},
dan \berkas{rekonsiliasi.md} yang memuat tabel perbandingan angka gudang lawan
angka sumber untuk ketiga query.
\end{luaran}

\begin{checkpoint}
Asisten menjalankan \berkas{sql/modul-02/check.sql} pada \perintah{nusamart\_dw}
dan memeriksa butir berikut:
\begin{ceklis}
  \item keempat tabel ada, dan \perintah{fakta\_penjualan} memiliki foreign key
        ke ketiga dimensi;
  \item \perintah{dim\_tanggal} memuat 3.653 baris beserta atribut kalender,
        bukan hanya kolom tanggal;
  \item \perintah{dim\_produk} dan \perintah{dim\_toko} memakai surrogate key
        yang terpisah dari natural key;
  \item tidak ada foreign key fakta yang \perintah{NULL};
  \item jumlah baris \perintah{fakta\_penjualan} sesuai dengan jumlah baris
        sumber setelah penyaringan status;
  \item angka ketiga star query sama persis dengan angka dari sumber;
  \item pernyataan grain tertulis sebagai komentar pada \berkas{ddl.sql}.
\end{ceklis}
\end{checkpoint}

\section{Latihan Mandiri di Kelas}

\latihan{Menambah satu atribut dimensi}

Tambahkan kolom \perintah{kelompok\_harga} pada \perintah{dim\_produk} yang
mengelompokkan produk menjadi \emph{ekonomis}, \emph{menengah}, dan
\emph{premium} berdasarkan harga jualnya. Lalu jawab:

\begin{enumerate}[leftmargin=1.4em]
  \item Tulis satu pertanyaan bisnis yang hanya dapat dijawab setelah kolom ini
        ada.
  \item Mengapa pengelompokan ini disimpan sebagai atribut dimensi, bukan
        dihitung ulang pada setiap query dengan \perintah{CASE WHEN}?
  \item Apa yang terjadi pada laporan tahun lalu bila harga sebuah produk naik
        sehingga kelompoknya berubah? Sebutkan modul mana yang menangani
        persoalan ini.
\end{enumerate}

\latihan{Menguji batas grain}

Jalankan query berikut, lalu jelaskan hasilnya.

\begin{lstlisting}[style=sqlgaya]
SELECT COUNT(*) AS baris_fakta,
       COUNT(DISTINCT (toko_key, nomor_struk)) AS jumlah_struk,
       ROUND(COUNT(*)::NUMERIC
             / COUNT(DISTINCT (toko_key, nomor_struk)), 2) AS produk_per_struk
FROM   fakta_penjualan;
\end{lstlisting}

Bila grain Anda ditetapkan pada tingkat struk dan bukan baris struk, kolom mana
pada query di atas yang menjadi mustahil dihitung? Kaitkan jawabannya dengan
kaidah ``grain halus dapat diringkas, grain kasar tidak dapat dipulihkan'' dari
Modul 1.

\section{Tugas Individu}

\subsection{Pekerjaan Wajib}
Transformasikan proses bisnis kedua dari model E/R sumber menjadi skema bintang
dan jelaskan setiap keputusan denormalisasi.

Pilih satu proses selain penjualan --- persediaan harian, cakupan promosi, atau
retur barang. Kumpulkan \berkas{skema-bintang-2.md} yang memuat:

\begin{enumerate}[leftmargin=1.4em]
  \item pernyataan grain dalam satu kalimat bisnis;
  \item DDL lengkap fakta dan dimensi barunya, tanpa perlu dijalankan;
  \item daftar dimensi yang \emph{dipakai bersama} dengan skema penjualan ---
        inilah calon \istilah{conformed dimension} yang Anda tandai pada bus
        matrix Modul 1;
  \item untuk setiap denormalisasi yang Anda ambil, satu kalimat yang menyebut
        apa yang dihemat dan apa yang dibayar.
\end{enumerate}

\subsection{Pertanyaan Analisis}

\begin{enumerate}[leftmargin=1.4em]
  \item Anda memuat \perintah{fakta\_penjualan} dengan \perintah{INNER JOIN} ke
        \perintah{dim\_produk}. Sebulan kemudian, sumber menambahkan produk baru
        yang belum masuk dimensi. Apa yang terjadi pada pemuatan berikutnya, dan
        mengapa kesalahan semacam ini sulit terdeteksi?
  \item Seorang rekan mengusulkan menyimpan kolom \perintah{margin\_persen} pada
        \perintah{fakta\_penjualan} agar tidak perlu dihitung berulang. Jelaskan
        mengapa usulan ini keliru, dan sebutkan apa yang seharusnya disimpan
        sebagai gantinya.
  \item Mengapa \perintah{dim\_tanggal} dibangkitkan sampai 2029, padahal data
        transaksi hanya sampai tahun berjalan? Apa akibatnya bila kalender hanya
        dibangkitkan sepanjang rentang data yang ada?
\end{enumerate}

\section{Format Pengumpulan dan Checklist}

Kumpulkan tiga skrip SQL yang dapat dijalankan berurutan. Sertakan hasil
rekonsiliasi dan pastikan dimensi tanggal tidak hanya berisi nilai tanggal.

Seluruh berkas diletakkan pada \berkas{modul-02-star-schema/}.

\begin{ceklis}
  \item \berkas{pre-lab.md}
  \item \berkas{ddl.sql} --- memuat pernyataan grain sebagai komentar
  \item \berkas{load.sql}
  \item \berkas{query.sql} --- memuat ketiga star query
  \item \berkas{rekonsiliasi.md}
  \item \berkas{skema-bintang-2.md}
  \item \berkas{jawaban-analisis.md}
  \item Ketiga skrip dapat dijalankan berurutan pada basis data kosong tanpa
        penyuntingan manual.
\end{ceklis}

\section{Troubleshooting}

\begin{table}[htbp]
\centering
\small
\caption{Masalah yang lazim muncul pada Modul 2 dan penanganannya}
\label{tab:m2-troubleshooting}
\begin{tabular}{@{}p{0.30\textwidth}>{\raggedright\arraybackslash}p{0.62\textwidth}@{}}
\toprule
\textbf{Gejala} & \textbf{Penanganan} \\
\midrule
\perintah{duplicate key value violates unique constraint} pada
  \perintah{dim\_produk} &
  Sumber memuat lebih dari satu baris untuk satu \perintah{produk\_id}, atau
  skrip pemuatan dijalankan dua kali. Periksa dengan
  \perintah{GROUP BY produk\_id HAVING COUNT(*) > 1}, lalu jadikan
  \berkas{load.sql} idempoten. \\
\perintah{insert or update violates foreign key constraint} pada fakta &
  Ada natural key pada sumber yang tidak memiliki pasangan di dimensi. Muat
  dimensi lebih dahulu, dan hitung berapa baris yang tidak menemukan pasangan. \\
Jumlah baris fakta lebih sedikit daripada sumber &
  \perintah{INNER JOIN} membuang baris tanpa pasangan dimensi. Ganti sementara
  dengan \perintah{LEFT JOIN} untuk menemukan barisnya, lalu catat sebabnya. \\
Angka star query lebih besar daripada sumber &
  Terjadi penggandaan baris akibat join ke dimensi yang tidak unik. Pastikan
  natural key setiap dimensi benar-benar \perintah{UNIQUE}. \\
Angka star query lebih kecil daripada sumber &
  Penyaringan status berbeda antara kedua sisi, atau baris fakta hilang saat
  pemuatan. Periksa \perintah{WHERE} kedua query. \\
Selisih kecil pada digit terakhir &
  Perbedaan pembulatan. Samakan presisi \perintah{NUMERIC} kedua sisi dan
  bandingkan dengan \perintah{ROUND(..., 2)}. \\
Nama hari dan bulan berbahasa Inggris &
  Setelan \perintah{lc\_time} basis data. Bukan kesalahan skrip; pilih satu
  bahasa lalu gunakan secara konsisten. \\
\bottomrule
\end{tabular}
\end{table}

\section{Ringkasan}

Modul ini mengubah satu kalimat grain menjadi empat tabel yang berjalan. Tiga
gagasan perlu dibawa ke Modul 3.

\emph{Pertama}, surrogate key bukan formalitas. Ia yang membuat satu produk
dapat memiliki beberapa versi tanpa memutus hubungan dengan tabel fakta --- dan
itulah tepatnya yang dikerjakan Modul 3. \emph{Kedua}, denormalisasi adalah
pertukaran yang harus dapat dijelaskan: ruang dan redundansi ditukar dengan
hilangnya join, dan pada gudang data pertukaran itu hampir selalu
menguntungkan. \emph{Ketiga}, angka yang tidak direkonsiliasi tidak dapat
dipercaya; selisih yang dijelaskan adalah temuan, sedangkan selisih yang
didiamkan adalah kegagalan.

Satu hal sengaja ditunda. \perintah{dim\_produk} yang Anda bangun hari ini hanya
menyimpan keadaan terkini: ketika harga sebuah produk berubah, versi lamanya
hilang dan laporan tahun lalu ikut berubah. Modul 3 memperbaiki persoalan ini
dengan SCD Type 2, dan struktur yang Anda buat sekarang sudah siap untuk itu.
